% Macros communes aux textes en vis-à-vis (Galates, Ézéchiel) :
% langue originale sur les pages de gauche, français sur les pages de droite,
% synchronisés verset par verset grâce au paquet paracol.
% À inclure dans le préambule : % Macros communes aux textes en vis-à-vis (Galates, Ézéchiel) :
% langue originale sur les pages de gauche, français sur les pages de droite,
% synchronisés verset par verset grâce au paquet paracol.
% À inclure dans le préambule : % Macros communes aux textes en vis-à-vis (Galates, Ézéchiel) :
% langue originale sur les pages de gauche, français sur les pages de droite,
% synchronisés verset par verset grâce au paquet paracol.
% À inclure dans le préambule : % Macros communes aux textes en vis-à-vis (Galates, Ézéchiel) :
% langue originale sur les pages de gauche, français sur les pages de droite,
% synchronisés verset par verset grâce au paquet paracol.
% À inclure dans le préambule : \input{../assets/versets.tex}
\usepackage{paracol}

%%%%% versets : \verset{langue originale}{français} %%%%%
% À l'intérieur de \verset, chaque ligne du fichier est une ligne du PDF.
% Les > en début de ligne donnent le niveau d'indentation,
% et un / après les > met toute la ligne en italique :
%   \verset
%     {Παῦλος ἀπόστολος
%      > οὐκ ἀπ’ ἀνθρώπων
%      >> διὰ Ἰησοῦ Χριστοῦ}
%     {Paul, apôtre,
%      > non de la part des hommes,
%      >> / par Jésus Christ}
% Les deux textes commencent à la même hauteur (synchronisation paracol).
% \horsverset{langue originale}{français} : même chose, sans numéro de verset ni indentation.
% \chapitre{n} : n'affiche rien, remet les versets à 1 et sert au numéro rouge du verset 1.
% \pagegauche : passe à la prochaine page paire, à utiliser juste avant \begin{paracol}[1]*{2}.
%
% Pour une langue écrite de droite à gauche, redéfinir l'environnement colonnegauche :
%   \renewenvironment{colonnegauche}{\begin{hebrew}}{\end{hebrew}}
\newenvironment{colonnegauche}{}{}

\def \pscolor{red} %couleur de la numérotation des versets
\def \pslabelsep{0.3em} %distance entre le numéro de verset et le texte
\def \psleftmargin{1.5em} %marge à gauche du texte des versets
\newlength{\indentunit}
\setlength{\indentunit}{0.5cm} %largeur d'un niveau d'indentation (un >)
\newlength{\versethang}
\setlength{\versethang}{1em} %retrait en plus quand une ligne trop longue se replie
\newcounter{chapitre}
\newcounter{verset}
% vrai après une section ou une sous-section : le verset suivant rappelle le chapitre
% (à activer avec \global\rappelchapitretrue dans le \titleformat du document)
\newif\ifrappelchapitre
\newcommand{\chapitre}[1]{\setcounter{chapitre}{#1}\setcounter{verset}{0}}
\newcommand{\pagegauche}{\clearpage\ifodd\value{page}\thispagestyle{empty}\null\clearpage\fi}
% numéro du verset, précédé du numéro de chapitre au verset 1
% et au premier verset qui suit une section ou une sous-section.
% Le tout occupe la marge \psleftmargin : le chapitre est calé sur le bord de la page (aligné
% avec les titres) et le verset contre le texte. De droite à gauche, ce placement s'inverse.
% (\textdir TLT sur chaque nombre : ses chiffres restent dans le bon sens)
\newcommand{\numverset}{\hbox to \psleftmargin{{\fontspec{Linux Libertine}\color{\pscolor}%
  \ifnum\value{verset}=1 \rappelchapitretrue\fi
  \ifrappelchapitre{\textdir TLT\arabic{chapitre}}\fi\hfil
  {\textdir TLT\small\textsuperscript{\arabic{verset}}}}\hspace{\pslabelsep}}}

% --- mécanique interne (pas besoin d'y toucher) ---
\makeatletter
\newcount\verset@niv
\newif\ifverset@premiere
% compte les > au début d'une ligne, puis applique le niveau
\def\verset@niveau{\verset@niv=0 \verset@compte}
\def\verset@compte{\futurelet\verset@suivant\verset@compteB}
\def\verset@compteB{\ifx\verset@suivant>\expandafter\verset@plus
  \else\expandafter\verset@applique\fi}
\def\verset@plus>{\advance\verset@niv\@ne \verset@compte}
\def\verset@applique{%
  \setlength{\leftskip}{\dimexpr\psleftmargin+\indentunit*\verset@niv\relax}%
  \hangindent\versethang \hangafter\@ne
  \ifverset@premiere \verset@premierefalse
    \noindent\llap{\numverset\hspace{\dimexpr\indentunit*\verset@niv\relax}}%
  \fi
  \verset@italique}
% saute les espaces, puis un / éventuel met la ligne en italique
\def\verset@italique{\futurelet\verset@suivant\verset@italiqueB}
\def\verset@italiqueB{%
  \ifx\verset@suivant\@sptoken \expandafter\verset@italiqueespace
  \else\ifx\verset@suivant/\expandafter\expandafter\expandafter\verset@italiquebarre
  \else \expandafter\expandafter\expandafter\ignorespaces\fi\fi}
\def\verset@italiqueespace{\afterassignment\verset@italique\let\verset@tmp= }
\def\verset@italiquebarre/{\itshape\ignorespaces}
% saute les fins de ligne et les espaces avant un argument
\def\verset@saute#1{\def\verset@apres{#1}\verset@sauteB}
\def\verset@sauteB{\futurelet\verset@suivant\verset@sauteC}
\def\verset@sauteC{%
  \ifx\verset@suivant\verset@nl \expandafter\verset@mange
  \else\ifx\verset@suivant\@sptoken \expandafter\expandafter\expandafter\verset@mangeespace
  \else \expandafter\expandafter\expandafter\verset@apres\fi\fi}
\def\verset@mange#1{\verset@sauteB}
\def\verset@mangeespace{\afterassignment\verset@sauteB\let\verset@tmp= }
% fin de ligne dans un verset : nouvelle ligne (l'italique s'arrête), puis lecture du niveau
\def\verset@nl{\par\upshape\verset@niveau}
\begingroup
\catcode`\^^M=\active%
\gdef\verset@lignes{\catcode`\^^M=\active\let^^M\verset@nl}%
\gdef\versetpar#1{{\setlength{\parskip}{0pt}\let^^M\verset@nl%
  \verset@premieretrue\verset@niveau#1\par}}%
\endgroup
\newcommand{\verset}{\begingroup\verset@lignes\verset@saute\verset@lireA}
\def\verset@lireA#1{\verset@saute{\verset@lireB{#1}}}
\def\verset@lireB#1#2{\endgroup\stepcounter{verset}%
  \switchcolumn[0]*\begin{colonnegauche}\versetpar{#1}\end{colonnegauche}%
  \switchcolumn[1]\versetpar{#2}%
  \global\rappelchapitrefalse}
\makeatother
\newcommand{\horsverset}[2]{\switchcolumn[0]*\begin{colonnegauche}{\setlength{\parskip}{0pt}#1\par}\end{colonnegauche}%
  \switchcolumn[1]{\setlength{\parskip}{0pt}#2\par}}

% compteurs communs aux deux colonnes (paracol les sépare par défaut)
\globalcounter{chapitre}
\globalcounter{verset}

\usepackage{paracol}

%%%%% versets : \verset{langue originale}{français} %%%%%
% À l'intérieur de \verset, chaque ligne du fichier est une ligne du PDF.
% Les > en début de ligne donnent le niveau d'indentation,
% et un / après les > met toute la ligne en italique :
%   \verset
%     {Παῦλος ἀπόστολος
%      > οὐκ ἀπ’ ἀνθρώπων
%      >> διὰ Ἰησοῦ Χριστοῦ}
%     {Paul, apôtre,
%      > non de la part des hommes,
%      >> / par Jésus Christ}
% Les deux textes commencent à la même hauteur (synchronisation paracol).
% \horsverset{langue originale}{français} : même chose, sans numéro de verset ni indentation.
% \chapitre{n} : n'affiche rien, remet les versets à 1 et sert au numéro rouge du verset 1.
% \pagegauche : passe à la prochaine page paire, à utiliser juste avant \begin{paracol}[1]*{2}.
%
% Pour une langue écrite de droite à gauche, redéfinir l'environnement colonnegauche :
%   \renewenvironment{colonnegauche}{\begin{hebrew}}{\end{hebrew}}
\newenvironment{colonnegauche}{}{}

\def \pscolor{red} %couleur de la numérotation des versets
\def \pslabelsep{0.3em} %distance entre le numéro de verset et le texte
\def \psleftmargin{1.5em} %marge à gauche du texte des versets
\newlength{\indentunit}
\setlength{\indentunit}{0.5cm} %largeur d'un niveau d'indentation (un >)
\newlength{\versethang}
\setlength{\versethang}{1em} %retrait en plus quand une ligne trop longue se replie
\newcounter{chapitre}
\newcounter{verset}
% vrai après une section ou une sous-section : le verset suivant rappelle le chapitre
% (à activer avec \global\rappelchapitretrue dans le \titleformat du document)
\newif\ifrappelchapitre
\newcommand{\chapitre}[1]{\setcounter{chapitre}{#1}\setcounter{verset}{0}}
\newcommand{\pagegauche}{\clearpage\ifodd\value{page}\thispagestyle{empty}\null\clearpage\fi}
% numéro du verset, précédé du numéro de chapitre au verset 1
% et au premier verset qui suit une section ou une sous-section.
% Le tout occupe la marge \psleftmargin : le chapitre est calé sur le bord de la page (aligné
% avec les titres) et le verset contre le texte. De droite à gauche, ce placement s'inverse.
% (\textdir TLT sur chaque nombre : ses chiffres restent dans le bon sens)
\newcommand{\numverset}{\hbox to \psleftmargin{{\fontspec{Linux Libertine}\color{\pscolor}%
  \ifnum\value{verset}=1 \rappelchapitretrue\fi
  \ifrappelchapitre{\textdir TLT\arabic{chapitre}}\fi\hfil
  {\textdir TLT\small\textsuperscript{\arabic{verset}}}}\hspace{\pslabelsep}}}

% --- mécanique interne (pas besoin d'y toucher) ---
\makeatletter
\newcount\verset@niv
\newif\ifverset@premiere
% compte les > au début d'une ligne, puis applique le niveau
\def\verset@niveau{\verset@niv=0 \verset@compte}
\def\verset@compte{\futurelet\verset@suivant\verset@compteB}
\def\verset@compteB{\ifx\verset@suivant>\expandafter\verset@plus
  \else\expandafter\verset@applique\fi}
\def\verset@plus>{\advance\verset@niv\@ne \verset@compte}
\def\verset@applique{%
  \setlength{\leftskip}{\dimexpr\psleftmargin+\indentunit*\verset@niv\relax}%
  \hangindent\versethang \hangafter\@ne
  \ifverset@premiere \verset@premierefalse
    \noindent\llap{\numverset\hspace{\dimexpr\indentunit*\verset@niv\relax}}%
  \fi
  \verset@italique}
% saute les espaces, puis un / éventuel met la ligne en italique
\def\verset@italique{\futurelet\verset@suivant\verset@italiqueB}
\def\verset@italiqueB{%
  \ifx\verset@suivant\@sptoken \expandafter\verset@italiqueespace
  \else\ifx\verset@suivant/\expandafter\expandafter\expandafter\verset@italiquebarre
  \else \expandafter\expandafter\expandafter\ignorespaces\fi\fi}
\def\verset@italiqueespace{\afterassignment\verset@italique\let\verset@tmp= }
\def\verset@italiquebarre/{\itshape\ignorespaces}
% saute les fins de ligne et les espaces avant un argument
\def\verset@saute#1{\def\verset@apres{#1}\verset@sauteB}
\def\verset@sauteB{\futurelet\verset@suivant\verset@sauteC}
\def\verset@sauteC{%
  \ifx\verset@suivant\verset@nl \expandafter\verset@mange
  \else\ifx\verset@suivant\@sptoken \expandafter\expandafter\expandafter\verset@mangeespace
  \else \expandafter\expandafter\expandafter\verset@apres\fi\fi}
\def\verset@mange#1{\verset@sauteB}
\def\verset@mangeespace{\afterassignment\verset@sauteB\let\verset@tmp= }
% fin de ligne dans un verset : nouvelle ligne (l'italique s'arrête), puis lecture du niveau
\def\verset@nl{\par\upshape\verset@niveau}
\begingroup
\catcode`\^^M=\active%
\gdef\verset@lignes{\catcode`\^^M=\active\let^^M\verset@nl}%
\gdef\versetpar#1{{\setlength{\parskip}{0pt}\let^^M\verset@nl%
  \verset@premieretrue\verset@niveau#1\par}}%
\endgroup
\newcommand{\verset}{\begingroup\verset@lignes\verset@saute\verset@lireA}
\def\verset@lireA#1{\verset@saute{\verset@lireB{#1}}}
\def\verset@lireB#1#2{\endgroup\stepcounter{verset}%
  \switchcolumn[0]*\begin{colonnegauche}\versetpar{#1}\end{colonnegauche}%
  \switchcolumn[1]\versetpar{#2}%
  \global\rappelchapitrefalse}
\makeatother
\newcommand{\horsverset}[2]{\switchcolumn[0]*\begin{colonnegauche}{\setlength{\parskip}{0pt}#1\par}\end{colonnegauche}%
  \switchcolumn[1]{\setlength{\parskip}{0pt}#2\par}}

% compteurs communs aux deux colonnes (paracol les sépare par défaut)
\globalcounter{chapitre}
\globalcounter{verset}

\usepackage{paracol}

%%%%% versets : \verset{langue originale}{français} %%%%%
% À l'intérieur de \verset, chaque ligne du fichier est une ligne du PDF.
% Les > en début de ligne donnent le niveau d'indentation,
% et un / après les > met toute la ligne en italique :
%   \verset
%     {Παῦλος ἀπόστολος
%      > οὐκ ἀπ’ ἀνθρώπων
%      >> διὰ Ἰησοῦ Χριστοῦ}
%     {Paul, apôtre,
%      > non de la part des hommes,
%      >> / par Jésus Christ}
% Les deux textes commencent à la même hauteur (synchronisation paracol).
% \horsverset{langue originale}{français} : même chose, sans numéro de verset ni indentation.
% \chapitre{n} : n'affiche rien, remet les versets à 1 et sert au numéro rouge du verset 1.
% \pagegauche : passe à la prochaine page paire, à utiliser juste avant \begin{paracol}[1]*{2}.
%
% Pour une langue écrite de droite à gauche, redéfinir l'environnement colonnegauche :
%   \renewenvironment{colonnegauche}{\begin{hebrew}}{\end{hebrew}}
\newenvironment{colonnegauche}{}{}

\def \pscolor{red} %couleur de la numérotation des versets
\def \pslabelsep{0.3em} %distance entre le numéro de verset et le texte
\def \psleftmargin{1.5em} %marge à gauche du texte des versets
\newlength{\indentunit}
\setlength{\indentunit}{0.5cm} %largeur d'un niveau d'indentation (un >)
\newlength{\versethang}
\setlength{\versethang}{1em} %retrait en plus quand une ligne trop longue se replie
\newcounter{chapitre}
\newcounter{verset}
% vrai après une section ou une sous-section : le verset suivant rappelle le chapitre
% (à activer avec \global\rappelchapitretrue dans le \titleformat du document)
\newif\ifrappelchapitre
\newcommand{\chapitre}[1]{\setcounter{chapitre}{#1}\setcounter{verset}{0}}
\newcommand{\pagegauche}{\clearpage\ifodd\value{page}\thispagestyle{empty}\null\clearpage\fi}
% numéro du verset, précédé du numéro de chapitre au verset 1
% et au premier verset qui suit une section ou une sous-section.
% Le tout occupe la marge \psleftmargin : le chapitre est calé sur le bord de la page (aligné
% avec les titres) et le verset contre le texte. De droite à gauche, ce placement s'inverse.
% (\textdir TLT sur chaque nombre : ses chiffres restent dans le bon sens)
\newcommand{\numverset}{\hbox to \psleftmargin{{\fontspec{Linux Libertine}\color{\pscolor}%
  \ifnum\value{verset}=1 \rappelchapitretrue\fi
  \ifrappelchapitre{\textdir TLT\arabic{chapitre}}\fi\hfil
  {\textdir TLT\small\textsuperscript{\arabic{verset}}}}\hspace{\pslabelsep}}}

% --- mécanique interne (pas besoin d'y toucher) ---
\makeatletter
\newcount\verset@niv
\newif\ifverset@premiere
% compte les > au début d'une ligne, puis applique le niveau
\def\verset@niveau{\verset@niv=0 \verset@compte}
\def\verset@compte{\futurelet\verset@suivant\verset@compteB}
\def\verset@compteB{\ifx\verset@suivant>\expandafter\verset@plus
  \else\expandafter\verset@applique\fi}
\def\verset@plus>{\advance\verset@niv\@ne \verset@compte}
\def\verset@applique{%
  \setlength{\leftskip}{\dimexpr\psleftmargin+\indentunit*\verset@niv\relax}%
  \hangindent\versethang \hangafter\@ne
  \ifverset@premiere \verset@premierefalse
    \noindent\llap{\numverset\hspace{\dimexpr\indentunit*\verset@niv\relax}}%
  \fi
  \verset@italique}
% saute les espaces, puis un / éventuel met la ligne en italique
\def\verset@italique{\futurelet\verset@suivant\verset@italiqueB}
\def\verset@italiqueB{%
  \ifx\verset@suivant\@sptoken \expandafter\verset@italiqueespace
  \else\ifx\verset@suivant/\expandafter\expandafter\expandafter\verset@italiquebarre
  \else \expandafter\expandafter\expandafter\ignorespaces\fi\fi}
\def\verset@italiqueespace{\afterassignment\verset@italique\let\verset@tmp= }
\def\verset@italiquebarre/{\itshape\ignorespaces}
% saute les fins de ligne et les espaces avant un argument
\def\verset@saute#1{\def\verset@apres{#1}\verset@sauteB}
\def\verset@sauteB{\futurelet\verset@suivant\verset@sauteC}
\def\verset@sauteC{%
  \ifx\verset@suivant\verset@nl \expandafter\verset@mange
  \else\ifx\verset@suivant\@sptoken \expandafter\expandafter\expandafter\verset@mangeespace
  \else \expandafter\expandafter\expandafter\verset@apres\fi\fi}
\def\verset@mange#1{\verset@sauteB}
\def\verset@mangeespace{\afterassignment\verset@sauteB\let\verset@tmp= }
% fin de ligne dans un verset : nouvelle ligne (l'italique s'arrête), puis lecture du niveau
\def\verset@nl{\par\upshape\verset@niveau}
\begingroup
\catcode`\^^M=\active%
\gdef\verset@lignes{\catcode`\^^M=\active\let^^M\verset@nl}%
\gdef\versetpar#1{{\setlength{\parskip}{0pt}\let^^M\verset@nl%
  \verset@premieretrue\verset@niveau#1\par}}%
\endgroup
\newcommand{\verset}{\begingroup\verset@lignes\verset@saute\verset@lireA}
\def\verset@lireA#1{\verset@saute{\verset@lireB{#1}}}
\def\verset@lireB#1#2{\endgroup\stepcounter{verset}%
  \switchcolumn[0]*\begin{colonnegauche}\versetpar{#1}\end{colonnegauche}%
  \switchcolumn[1]\versetpar{#2}%
  \global\rappelchapitrefalse}
\makeatother
\newcommand{\horsverset}[2]{\switchcolumn[0]*\begin{colonnegauche}{\setlength{\parskip}{0pt}#1\par}\end{colonnegauche}%
  \switchcolumn[1]{\setlength{\parskip}{0pt}#2\par}}

% compteurs communs aux deux colonnes (paracol les sépare par défaut)
\globalcounter{chapitre}
\globalcounter{verset}

\usepackage{paracol}

%%%%% versets : \verset{langue originale}{français} %%%%%
% À l'intérieur de \verset, chaque ligne du fichier est une ligne du PDF.
% Les > en début de ligne donnent le niveau d'indentation,
% et un / après les > met toute la ligne en italique :
%   \verset
%     {Παῦλος ἀπόστολος
%      > οὐκ ἀπ’ ἀνθρώπων
%      >> διὰ Ἰησοῦ Χριστοῦ}
%     {Paul, apôtre,
%      > non de la part des hommes,
%      >> / par Jésus Christ}
% Les deux textes commencent à la même hauteur (synchronisation paracol).
% \horsverset{langue originale}{français} : même chose, sans numéro de verset ni indentation.
% \chapitre{n} : n'affiche rien, remet les versets à 1 et sert au numéro rouge du verset 1.
% \pagegauche : passe à la prochaine page paire, à utiliser juste avant \begin{paracol}[1]*{2}.
%
% Pour une langue écrite de droite à gauche, redéfinir l'environnement colonnegauche :
%   \renewenvironment{colonnegauche}{\begin{hebrew}}{\end{hebrew}}
\newenvironment{colonnegauche}{}{}

\def \pscolor{red} %couleur de la numérotation des versets
\def \pslabelsep{0.3em} %distance entre le numéro de verset et le texte
\def \psleftmargin{1.5em} %marge à gauche du texte des versets
\newlength{\indentunit}
\setlength{\indentunit}{0.5cm} %largeur d'un niveau d'indentation (un >)
\newlength{\versethang}
\setlength{\versethang}{1em} %retrait en plus quand une ligne trop longue se replie
\newcounter{chapitre}
\newcounter{verset}
% vrai après une section ou une sous-section : le verset suivant rappelle le chapitre
% (à activer avec \global\rappelchapitretrue dans le \titleformat du document)
\newif\ifrappelchapitre
\newcommand{\chapitre}[1]{\setcounter{chapitre}{#1}\setcounter{verset}{0}}
\newcommand{\pagegauche}{\clearpage\ifodd\value{page}\thispagestyle{empty}\null\clearpage\fi}
% numéro du verset, précédé du numéro de chapitre au verset 1
% et au premier verset qui suit une section ou une sous-section.
% Le tout occupe la marge \psleftmargin : le chapitre est calé sur le bord de la page (aligné
% avec les titres) et le verset contre le texte. De droite à gauche, ce placement s'inverse.
% (\textdir TLT sur chaque nombre : ses chiffres restent dans le bon sens)
\newcommand{\numverset}{\hbox to \psleftmargin{{\fontspec{Linux Libertine}\color{\pscolor}%
  \ifnum\value{verset}=1 \rappelchapitretrue\fi
  \ifrappelchapitre{\textdir TLT\arabic{chapitre}}\fi\hfil
  {\textdir TLT\small\textsuperscript{\arabic{verset}}}}\hspace{\pslabelsep}}}

% --- mécanique interne (pas besoin d'y toucher) ---
\makeatletter
\newcount\verset@niv
\newif\ifverset@premiere
% compte les > au début d'une ligne, puis applique le niveau
\def\verset@niveau{\verset@niv=0 \verset@compte}
\def\verset@compte{\futurelet\verset@suivant\verset@compteB}
\def\verset@compteB{\ifx\verset@suivant>\expandafter\verset@plus
  \else\expandafter\verset@applique\fi}
\def\verset@plus>{\advance\verset@niv\@ne \verset@compte}
\def\verset@applique{%
  \setlength{\leftskip}{\dimexpr\psleftmargin+\indentunit*\verset@niv\relax}%
  \hangindent\versethang \hangafter\@ne
  \ifverset@premiere \verset@premierefalse
    \noindent\llap{\numverset\hspace{\dimexpr\indentunit*\verset@niv\relax}}%
  \fi
  \verset@italique}
% saute les espaces, puis un / éventuel met la ligne en italique
\def\verset@italique{\futurelet\verset@suivant\verset@italiqueB}
\def\verset@italiqueB{%
  \ifx\verset@suivant\@sptoken \expandafter\verset@italiqueespace
  \else\ifx\verset@suivant/\expandafter\expandafter\expandafter\verset@italiquebarre
  \else \expandafter\expandafter\expandafter\ignorespaces\fi\fi}
\def\verset@italiqueespace{\afterassignment\verset@italique\let\verset@tmp= }
\def\verset@italiquebarre/{\itshape\ignorespaces}
% saute les fins de ligne et les espaces avant un argument
\def\verset@saute#1{\def\verset@apres{#1}\verset@sauteB}
\def\verset@sauteB{\futurelet\verset@suivant\verset@sauteC}
\def\verset@sauteC{%
  \ifx\verset@suivant\verset@nl \expandafter\verset@mange
  \else\ifx\verset@suivant\@sptoken \expandafter\expandafter\expandafter\verset@mangeespace
  \else \expandafter\expandafter\expandafter\verset@apres\fi\fi}
\def\verset@mange#1{\verset@sauteB}
\def\verset@mangeespace{\afterassignment\verset@sauteB\let\verset@tmp= }
% fin de ligne dans un verset : nouvelle ligne (l'italique s'arrête), puis lecture du niveau
\def\verset@nl{\par\upshape\verset@niveau}
\begingroup
\catcode`\^^M=\active%
\gdef\verset@lignes{\catcode`\^^M=\active\let^^M\verset@nl}%
\gdef\versetpar#1{{\setlength{\parskip}{0pt}\let^^M\verset@nl%
  \verset@premieretrue\verset@niveau#1\par}}%
\endgroup
\newcommand{\verset}{\begingroup\verset@lignes\verset@saute\verset@lireA}
\def\verset@lireA#1{\verset@saute{\verset@lireB{#1}}}
\def\verset@lireB#1#2{\endgroup\stepcounter{verset}%
  \switchcolumn[0]*\begin{colonnegauche}\versetpar{#1}\end{colonnegauche}%
  \switchcolumn[1]\versetpar{#2}%
  \global\rappelchapitrefalse}
\makeatother
\newcommand{\horsverset}[2]{\switchcolumn[0]*\begin{colonnegauche}{\setlength{\parskip}{0pt}#1\par}\end{colonnegauche}%
  \switchcolumn[1]{\setlength{\parskip}{0pt}#2\par}}

% compteurs communs aux deux colonnes (paracol les sépare par défaut)
\globalcounter{chapitre}
\globalcounter{verset}
